\documentclass[12pt]{article}
\usepackage[utf8]{inputenc}
\usepackage[margin=1in]{geometry}
\usepackage{indentfirst}
\usepackage[osf]{mathpazo}

\date{} %date not required
%\title{Competing for Quality: Coordination Failure in Dairy Genetics Markets} 

%\author{
%    Victor Funes-Leal \\ % Author 1
%    Department of Agricultural Economics and Agribusiness,\\
%    University of Arkansas \\
%    \texttt{vf006@uark.edu}
%    \and % Separator for multiple authors
%    Jared Hutchins\\ % Author 2
%    Department of Agricultural and Consumer Economics,\\
%    University of Illinois Urbana Champaign \\
%    \texttt{jhtchns2@illinois.edu}
%}


\begin{document}
%\maketitle

Dear editor,

Thank you for the opportunity to resubmit our article, which he have now titled “Competing for Quality: Coordination Failure in Dairy Genetics Market.” We have reviewed all three reviewers' comments and made several revisions to address them. The major changes made to the paper are summarized below:

\begin{itemize}
    \item Reviewers 1 and 3 questioned our framing of the problem, which caused increased inbreeding rates among US Holstein cattle, as an externality. We have included a new section (Section 2) that addresses their questions using a mathematical model that describes the economic process underlying the effects we observe in the data. In that model, we have stepped away from framing the issue as an externality. Instead, we frame the issue as a coordination problem in which competition incentivizes farmers to inbreed, leading them to compete without considering the positive benefits of outcrossing for long-run industry growth.
    \item Reviewer 1 raised several questions about our empirical strategy, among them, pointing out that we are estimating an effect on the treated (high-quality sire progeny), not the entire industry, as we mistakenly pointed out. All remaining questions have been addressed in the corresponding section (Section 4).
    \item All three reviewers requested significant changes to our cost-benefit analysis; they noted that we did not include a Net Present Value calculation because we did not discount the values. Reviewer 3 pointed out that we cannot conduct a proper cost-benefit analysis, and we have included discounting and added an appendix that explains our assumptions and the scope of our calculations.
    \item Reviewer 2 asked for a formal model that explains animal breeders' choices between inbreeding versus outcrossing their bulls. We have a new section (Section 2.1) that formalizes the breeder's choices and their consequences.
    \item Reviewer 3 was concerned that our paper was not a good fit for an economics journal because of the perception that our paper was calculating inbreeding rates. We have added a succinct explanation of the concept, explaining that it is calculated and reported by the NAAB. Finally, we explain that our interest is to identify the economic forces that lead to an undesirable outcome, which ultimately results in a series of measurable negative financial impacts for the dairy industry.
    \item Reviewers 1 and 2 show that we did not make a proper distinction between dairy farmers and animal breeders, the former being those who demand bull genetics, while the latter are those who supply such goods.
    \item Reviewer 3 mentioned that we did not properly explain why the number of bulls in the market increased so rapidly. We added a line showing that this is the result of genetics companies changing their marketing decisions by releasing younger bulls into the market.
    \item Reviewer 2 requested a clarification about the decision process of dairy farmers: why are some bulls so popular while others are not? We explain that pedigree is a key deciding factor, especially among top sires whose genetic traits are in the upper tail of the distribution, making one-to-one comparisons between bulls extremely difficult.
    \end{itemize}

Thank you for considering our work for publication in the \textbf{American Journal of Agricultural and Economics}.
\\

Sincerely,

Victor Funes-Leal, Ph.D.

Jared Hutchins, Ph.D.

\end{document}